% ──────────────────────────────────────────────────────────────────────────────
% INTRODUCTION DRAFT
% ──────────────────────────────────────────────────────────────────────────────

\section{Introduction}

The role of machine learning models in high-stakes decision-making continues to
expand, shaping outcomes in finance, employment, insurance, and healthcare.
When these same models perform systematically worse for particular demographic
or protected groups, the resulting errors can translate into unequal access to
opportunities, heightened financial risk, and measurable downstream harms.
These disparities arise from a combination of data, measurement, representation,
and deployment dynamics, and can persist even when overall accuracy appears
strong \cite{SureshGuttag2021SourcesOfHarm}.

A core technical challenge is that standard empirical risk minimization (ERM),
which serves as the foundation for the vast majority of modern supervised
learning systems, optimizes average performance and can obscure poor outcomes
on minority or historically disadvantaged subpopulations
\cite{Mehrabi2021SurveyBiasFairness}.

In practice, these failures often become visible during audits or
post-deployment monitoring, motivating standardized reporting and evaluation
tools (e.g., model documentation and subgroup evaluations)
\cite{Mitchell2019ModelCards,Gebru2018Datasheets}.
Empirical evidence in deployed settings further reinforces that performance can
vary substantially across demographic intersections even for commercially used
systems, underscoring the importance of subgroup-aware evaluation rather than
relying on headline performance alone \cite{BuolamwiniGebru2018GenderShades}.

A substantial body of work addresses fairness by modifying the training
objective of a classifier.
Loss regularization, constraint-based optimization, adversarial debiasing, and
distributionally robust optimization incorporate fairness criteria directly into
the learning process \cite{sagawa2020distributionally}.
These methods explicitly penalize disparities in error rates or optimize
worst-group risk during training.

While such approaches can improve specific fairness metrics, they are
inherently model-dependent.
Fairness constraints must typically be derived for particular architectures or
loss functions, and the resulting fairness--utility trade-offs are often
sensitive to hyperparameters, constraint weights, and training dynamics
\cite{Mehrabi2021SurveyBiasFairness}.

Moreover, fairness improvements achieved through model-level interventions do
not necessarily transfer across downstream tasks.
In practical settings where multiple models are trained on shared datasets,
potentially with heterogeneous architectures, embedding fairness constraints
into each model pipeline increases engineering complexity and deployment
friction.
These considerations motivate exploration of fairness interventions that are
less tightly coupled to specific classifier objectives.

Data-level approaches attempt to mitigate bias prior to model training.
Techniques such as reweighting, resampling, representation learning, and
feature modification aim to adjust the empirical training distribution so that
standard learning algorithms produce more equitable outcomes
\cite{kamiran2012data}.
Synthetic data methods extend this paradigm by generating additional samples
for underrepresented or disadvantaged groups.

A particularly challenging regime — which we term positive-class outcome
scarcity — arises when historical discrimination has directly suppressed
favourable outcomes for disadvantaged groups, leaving training data with
structurally few positive-class examples for those groups.
This pattern is directly observable in benchmark fairness datasets: in the
UCI Adult Census Income dataset, the positive-class rate (income $>$\$50K)
is approximately 30\% for male and 11\% for female individuals
\cite{goldin2000orchestrating}, reflecting documented employment and wage
barriers; analogous disparities have been identified in medical AI systems
trained on sex-imbalanced clinical records \cite{cirillo2020sex} and in
credit markets where access restrictions have produced corpora with
structurally fewer positive outcomes for women \cite{bartlett2022consumer}.
In this regime, reweighting and redistribution-based methods face a
structural constraint: with only a handful of minority-positive examples
available, there is insufficient signal to construct reliable
group-conditional gradients or transport plans, and oversampling amplifies
noise rather than true signal.
This motivates a generative approach — synthesising new minority-positive
samples can improve representation without requiring a sufficient supply of
existing positive examples to reweight from.

However, many existing data-level interventions rely on static heuristics or
distribution-matching objectives.
Synthetic data generation methods are often designed to replicate the original
data distribution or to enforce fixed fairness constraints during generation
\cite{Xu2019CTGAN}.
While such methods may improve representation balance, they are not explicitly
optimized for downstream worst-group performance.

Crucially, fairness metrics, particularly loss-based notions such as
worst-group risk, are defined in terms of model outcomes rather than purely
distributional properties.
Without feedback from downstream validation behaviour, data-level adjustments
may fail to target the specific regions of feature space responsible for
subgroup degradation.
This gap suggests a need for adaptive, outcome-aware data interventions.

In this work, we formulate fairness improvement as an outcome-driven synthetic
data generation problem.
Instead of modifying the classifier objective, we intervene at the data level
by learning to generate synthetic samples that reshape the effective training
distribution.
Generation is explicitly guided by feedback from downstream group-conditional
validation performance.

We model synthetic data generation as a sequential decision process and
optimize a generation policy using reinforcement learning.
The policy receives signals derived from worst-group validation loss, enabling
it to adaptively target regions of the feature space associated with
disadvantaged-group error concentration.
By coupling generation to outcome-level objectives, the framework directly
optimizes for worst-case group performance rather than distributional
similarity alone.

This formulation offers two key advantages.
First, fairness optimization remains largely model-agnostic: the generation
policy interacts with downstream models through outcome metrics rather than
architecture-specific constraints.
Second, the adaptive feedback loop allows the synthetic data mechanism to
focus on empirically observed failure modes, aligning data augmentation with
measurable subgroup risk.

The contributions of this work are as follows:
\begin{itemize}
    \item We formulate fairness-aware synthetic data generation as an
    outcome-driven optimization problem in which generation is guided by
    downstream worst-group validation risk.

    \item We introduce a reinforcement learning framework that adaptively
    reshapes the training distribution through targeted synthetic sample
    generation, without modifying the classifier objective.

    \item We demonstrate that the framework maintains low worst-group disparity
    under severe positive-class scarcity ($p_+ \leq 10\%$), a regime in which
    reweighting-based baselines degrade substantially, while preserving
    competitive predictive performance across multiple datasets and protected
    attributes.
\end{itemize}
